% SPDX-License-Identifier: CC0-1.0
% A starting document from the extension's templates, dedicated to the public domain (CC0 1.0). Replace the
% text and keep the structure you need; the preview updates as you type.
\documentclass[11pt]{article}
\usepackage[T1]{fontenc}
\usepackage[utf8]{inputenc}
\usepackage{amsmath,amssymb}
\usepackage{graphicx}
\usepackage{booktabs}
\usepackage[numbers]{natbib}
\usepackage[hidelinks]{hyperref}

\title{A short, informative title}
\author{First Author \and Second Author}
\date{\today}

\begin{document}
\maketitle

\begin{abstract}
One paragraph: the question, what you did, what you found, and why it matters.
\end{abstract}

\section{Introduction}
\label{sec:intro}
State the problem and why it matters. Cite sources as \cite{callaway1959} or \cite{slack1973,ziman1960}, and
refer to Section~\ref{sec:method}, Equation~\eqref{eq:model} and Table~\ref{tab:results}.

\section{Method}
\label{sec:method}
The kinetic model gives the thermal conductivity as
\begin{equation}
  \kappa = \frac{1}{3} C v \ell .
  \label{eq:model}
\end{equation}

\section{Results}
\begin{table}[htbp]
  \centering
  \caption{A table with booktabs rules.}
  \label{tab:results}
  \begin{tabular}{lrr}
    \toprule
    Sample & Value & Error \\
    \midrule
    A & 1.0 & 0.1 \\
    B & 2.0 & 0.2 \\
    \bottomrule
  \end{tabular}
\end{table}

\section{Conclusion}
What the reader should take away.

\bibliographystyle{plainnat}
\bibliography{references}
\end{document}
